\documentclass[12pt, a4paper]{article}

% --- PACOTES --- % 
\usepackage[utf8]{inputenc}
\usepackage[brazil]{babel}
\usepackage{geometry}
\usepackage{amsmath}
\usepackage{graphicx}
\usepackage{hyperref}
\usepackage{array}
\usepackage{longtable}
\usepackage{lmodern}
\usepackage{xcolor}

% --- CONFIGURAÇÕES --- %
\geometry{a4paper, top=3cm, bottom=3cm, left=3cm, right=3cm}
\hypersetup{
    colorlinks=true,
    linkcolor=blue,
    filecolor=magenta,      
    urlcolor=cyan,
}
\setlength{\parindent}{1.25cm}
\setlength{\parskip}{0.2cm}
\renewcommand{\baselinestretch}{1.5}

% --- TÍTULO --- %
\title{
    \textbf{Proposta de Trabalho de Conclusão de Curso} \\
    \vspace{1cm}
    \large Avaliação Abrangente de Modelos de Aprendizado de Máquina para a Classificação e Interpretabilidade de Materiais Supercondutores: \\
    Uma Extensão Metodológica do Trabalho Inicial
}
\author{Lucas Martin De Lucca \\ }
\date{\today}

% --- DOCUMENTO --- %
\begin{document}

\maketitle
\newpage

\tableofcontents
\newpage

\section{Introdução}

A descoberta de novos materiais com propriedades desejadas, como a supercondutividade em altas temperaturas, representa um dos maiores desafios da ciência de materiais contemporânea. A abordagem tradicional, baseada em experimentação e intuição científica, embora fundamental, é frequentemente lenta e custosa. O aprendizado de máquina (ML) emergiu como uma ferramenta poderosa para acelerar esse processo, permitindo a predição de propriedades a partir de características estruturais e eletrônicas calculadas computacionalmente.

Em trabalho anterior, desenvolvido no contexto da disciplina de Física Computacional da UFABC, demonstrei a viabilidade de classificar materiais como supercondutores ou não-supercondutores usando dados de Teoria do Funcional da Densidade (DFT). O estudo comparou três modelos de aprendizado de máquina (Random Forest, SVM e XGBoost) aplicados a um dataset de 2.474 materiais com 132 características de estrutura eletrônica, concluindo que o XGBoost obteve o melhor desempenho com F1-score de 0,7304.

Esta proposta de Trabalho de Conclusão de Curso (TCC) visa expandir significativamente o trabalho inicial, construindo um framework experimental robusto para avaliar e comparar sistematicamente um leque mais amplo de técnicas de ML, métodos de tratamento e seleção de características, e abordagens avançadas de interpretabilidade.

\section{Justificativa}

A expansão proposta é justificada por um aprofundamento da analise inicial. Primeiro, a inclusão de modelos de ponta como LightGBM e redes neurais, além de abordagens focadas em interpretabilidade como SISSO (Sure Independence Screening and Sparsifying Operator), permitirá uma avaliação mais completa do estado da arte para este problema específico. O trabalho original limitou-se a três algoritmos, deixando lacunas importantes na compreensão de qual abordagem é verdadeiramente ótima para este domínio.

Segundo, a análise comparativa de diferentes métodos de tratamento e seleção de características pode levar a um aumento significativo no desempenho preditivo. O trabalho inicial utilizou apenas imputação pela mediana e StandardScaler, sem explorar técnicas mais sofisticadas como IterativeImputer (MICE) ou QuantileTransformer, que podem ser mais adequadas para dados com alta dimensionalidade e valores ausentes.

Terceiro, e mais importante do ponto de vista científico, a aplicação de técnicas avançadas de interpretabilidade como SHAP (SHapley Additive exPlanations) e LIME (Local Interpretable Model-agnostic Explanations), em conjunto com o SISSO, pode fornecer insights valiosos sobre os mecanismos físicos que governam a supercondutividade. Esta análise transcende a simples predição, contribuindo para o conhecimento científico fundamental na área de ciência de materiais.

\section{Objetivos}

\subsection{Objetivo Geral}

Desenvolver e avaliar um framework abrangente de aprendizado de máquina para a classificação de materiais supercondutores, expandindo metodologicamente o trabalho inicial com foco na otimização do desempenho preditivo e na extração de conhecimento interpretável a partir dos modelos.

\subsection{Objetivos Específicos}

\begin{enumerate}
    \item Expandir o conjunto de modelos avaliados de 3 para 7 algoritmos distintos, incluindo modelos de gradient boosting (XGBoost, LightGBM), ensembles (Random Forest), SVM, modelos lineares regularizados (LASSO), redes neurais (MLP) e o SISSO.
    \item Comparar sistematicamente o impacto de 2 métodos avançados para tratamento de características (imputação iterativa e transformação por quantis) em relação às técnicas baseline utilizadas no trabalho inicial.
    \item Avaliar a eficácia de 3 estratégias diferentes para seleção de características: filtro (SelectKBest - baseline), wrapper (RFE) e embutido (LASSO), superando a abordagem única do trabalho original.
    \item Implementar e comparar 3 técnicas de interpretabilidade (LIME, SHAP e análise de descritores do SISSO) para explicar as predições dos modelos e extrair insights físicos sobre a supercondutividade.
    \item Identificar a combinação ótima de pré-processamento, seleção de características e modelo que maximize o desempenho da classificação, superando o F1-score de 0,7304 obtido anteriormente.
    \item Documentar todo o processo em um trabalho monográfico rigoroso, apresentando os resultados de forma clara, reprodutível e com análise estatística apropriada.
\end{enumerate}

\section{Metodologia Proposta}

A metodologia será estruturada como uma extensão sistemática do trabalho inicial, mantendo a base sólida já estabelecida e expandindo cada componente de forma controlada e científica.

\subsection{Dataset e Divisão dos Dados}
Utilizarei o mesmo dataset do trabalho inicial (2.474 materiais com 132 características de estrutura eletrônica do repositório SciDB), garantindo comparabilidade direta dos resultados. A divisão será aprimorada para incluir um conjunto de validação explícito: treino (70\%), validação (15\%) e teste (15\%), de forma estratificada para manter a proporção de classes. A validação cruzada de 5 folds será aplicada no conjunto de treino para otimização de hiperparâmetros e avaliação da estabilidade dos modelos.

\subsection{Tratamento Aprimorado de Características}
Expandindo as técnicas básicas do trabalho inicial, serão comparadas duas abordagens para cada etapa:

\textbf{Imputação de Dados Ausentes:}
\begin{itemize}
    \item \textit{Baseline:} Imputação pela mediana (método original)
    \item \textit{Avançado:} IterativeImputer (MICE) - método iterativo que modela cada característica como função das outras
\end{itemize}

\textbf{Escalonamento de Características:}
\begin{itemize}
    \item \textit{Baseline:} StandardScaler (método original)
    \item \textit{Avançado:} QuantileTransformer - mapeia para distribuição uniforme/normal, mais robusto a outliers
\end{itemize}

\subsection{Seleção Sistemática de Características}
Superando a abordagem única do trabalho inicial (SelectKBest), três métodos serão avaliados para selecionar as 50 características mais relevantes:

\begin{itemize}
    \item \textbf{Filtro (Baseline):} SelectKBest com teste F-ANOVA (método original)
    \item \textbf{Wrapper:} Recursive Feature Elimination (RFE) com Random Forest como estimador
    \item \textbf{Embutido:} Seleção baseada em regularização L1 (LASSO)
\end{itemize}

\subsection{Modelos de Aprendizado de Máquina}
Expandindo significativamente o conjunto original de 3 modelos, sete algoritmos serão treinados e avaliados:

\begin{longtable}{|p{3cm}|p{9cm}|}
\hline
\textbf{Modelo} & \textbf{Descrição e Justificativa} \\
\hline
\endfirsthead
\hline
\textbf{Modelo} & \textbf{Descrição e Justificativa} \\
\hline
\endhead
\hline
\endfoot
\endlastfoot
Random Forest & Baseline de ensemble do trabalho original, para comparação direta. \\
\hline
XGBoost & Melhor modelo do trabalho original (F1=0,7304), servirá como benchmark a ser superado. \\
\hline
SVM & Modelo de kernel do trabalho original, para análise de fronteiras de decisão não-lineares. \\
\hline
LASSO & Modelo linear regularizado que realiza seleção de características simultaneamente à classificação. \\
\hline
LightGBM & Alternativa moderna ao XGBoost, conhecida pela eficiência computacional e performance competitiva. \\
\hline
Rede Neural (MLP) & Para capturar relações complexas e não-lineares que modelos baseados em árvores podem não detectar. \\
\hline
SISSO & Foco em interpretabilidade física, buscando descritores na forma de equações analíticas interpretáveis. \\
\hline
\end{longtable}

Todos os hiperparâmetros serão otimizados por busca aleatória (Random Search) com validação cruzada, seguindo a metodologia do trabalho original mas com espaço de busca expandido.

\subsection{Avaliação Estatística}
O desempenho será avaliado usando o mesmo conjunto de métricas do trabalho original (Acurácia, Precisão, Recall, F1-Score, ROC-AUC) acrescido de Precision-Recall AUC. A significância estatística das diferenças de desempenho entre os modelos será verificada com testes estatísticos apropriados (teste t pareado, teste de Friedman para comparações múltiplas).

\subsection{Análise de Interpretabilidade}
Esta é a principal inovação em relação ao trabalho original. Três abordagens complementares serão aplicadas:

\begin{itemize}
    \item \textbf{Importância de Características Globais:} Extraída de modelos como Random Forest, XGBoost e LASSO, expandindo a análise superficial do trabalho original.
    \item \textbf{LIME:} Para explicações locais de predições individuais, permitindo entender por que um material específico foi classificado como supercondutor.
    \item \textbf{SHAP:} Para explicações locais e globais consistentes baseadas em teoria dos jogos, oferecendo uma visão unificada do comportamento dos modelos.
\end{itemize}

\section{Plano de Trabalho e Cronograma}

O trabalho será desenvolvido ao longo de 6 meses, seguindo um cronograma estruturado:

\begin{center}
\begin{tabular}{|l|c|c|c|c|c|c|}
\hline
\textbf{Atividade} & \textbf{Mês 1} & \textbf{2} & \textbf{3} & \textbf{4} & \textbf{5} & \textbf{6} \\
\hline
Revisão Bibliográfica Expandida & X & X & & & & \\
\hline
Preparação do Ambiente Computacional & X & X & & & & \\
\hline
Implementação da Pipeline Expandida & & X & X & & & \\
\hline
Execução dos Experimentos & & & X & X & & \\
\hline
Análise de Interpretabilidade & & & & X & & \\
\hline
Análise Estatística dos Resultados & & & & X & X & \\
\hline
Redação da Monografia & & & & & X & X \\
\hline
Revisão e Preparação da Defesa & & & & & & X \\
\hline
\end{tabular}
\end{center}

\section{Resultados Esperados}

Com base na experiência do trabalho inicial e nas melhorias metodológicas propostas, espero alcançar os seguintes resultados:

\begin{enumerate}
    \item \textbf{Melhoria de Performance:} Superar o F1-score de 0,7304 obtido no trabalho original através da otimização da pipeline completa.
    \item \textbf{Ranking Robusto de Características:} Identificar as características de estrutura eletrônica mais importantes para a supercondutividade, validado por múltiplos modelos e técnicas de interpretabilidade.
    \item \textbf{Análise Comparativa Definitiva:} Estabelecer qual combinação de pré-processamento, seleção de características e modelo é ótima para este problema específico.
    \item \textbf{Insights Físicos:} Descobrir novas relações entre propriedades eletrônicas e supercondutividade através das análises com SISSO, SHAP e LIME.
    \item \textbf{Framework Reprodutível:} Documentar uma metodologia completa e reprodutível que possa ser aplicada a outros problemas de descoberta de materiais.
\end{enumerate}

\section{Contribuições Esperadas}

Este trabalho representará uma contribuição significativa em três níveis:

\textbf{Metodológico:} Estabelecimento de um protocolo robusto para aplicação de ML em descoberta de materiais supercondutores, superando as limitações identificadas no trabalho inicial.

\textbf{Técnico:} Identificação da melhor pipeline de ML para este problema específico, com potencial de aplicação em outros domínios de ciência de materiais.

\textbf{Científico:} Extração de conhecimento interpretável sobre os mecanismos que governam a supercondutividade, contribuindo para o avanço do conhecimento fundamental na área.


\end{document}

